\begin{ouverture}[calcul]{11}{Proportionnalité et pourcentages}{Partie 2 · Algèbre et grandeurs}
{Coefficient · Retour à l’unité · Échelle · Vitesse moyenne}
\objectif{Reconnaître et représenter une proportionnalité.}
\objectif{Utiliser proportions et pourcentages.}
\objectif{Résoudre des problèmes de prix, recettes et échelles.}
\prerequis{Quotients et fractions.}
\end{ouverture}
\pagestyle{otMain}\essentiel\fontsize{11}{15}\selectfont
\phantomsection\label{cours-11-1}
\essSection{1}{Un même multiplicateur}
\begin{definition}Deux grandeurs sont proportionnelles si les valeurs de l’une s’obtiennent en multipliant les valeurs correspondantes de l’autre par un même nombre : le coefficient de proportionnalité.\end{definition}
\begin{center}\renewcommand{\arraystretch}{1.35}\begin{tabular}{c|ccc}\toprule
Masse (kg) & 1 & 2 & 3,5\\\midrule
Prix (euros) & 2,4 & 4,8 & 8,4\\
\bottomrule\end{tabular}\end{center}
\begin{savoirfaire}[Plusieurs procédures]On peut revenir à l’unité, additionner deux colonnes ou multiplier une colonne. Pour six personnes à partir d’une recette pour quatre, multiplier les quantités par $6/4=1{,}5$.\end{savoirfaire}
\label{fig-cours-11-1-cours-proportion}\figureGeo{cours-proportion}

\begin{attention}Un supplément fixe, un tarif par paliers ou une remise selon la quantité peut rompre la proportionnalité. Il faut vérifier le modèle avant de choisir la méthode.\end{attention}
\clearpage
\phantomsection\label{cours-11-2}
\essSection{2}{Pourcentages et changements d’échelle}
\begin{definition}$p\,\%$ signifie $p/100$. Les $15\,\%$ de 80 valent $(80:100)\times15=12$. Pour exprimer une part de 7 sur 20 en pourcentage, calculer $7/20=35/100=35\,\%$.\end{definition}
\begin{savoirfaire}[Échelles]À l’échelle $1:200$, une longueur de 1 cm représente 200 cm réels. Utiliser la même unité pour le dessin et la réalité avant de multiplier ou diviser.\end{savoirfaire}
\begin{exemple}À vitesse constante de 60 km/h, 2,5 heures correspondent à $60\times2{,}5=150$ km. Pour un déplacement quelconque, la vitesse moyenne est distance totale divisée par durée totale. 30 minutes correspondent à 0,5 heure.\end{exemple}
\begin{propriete}Dans un repère, une proportionnalité donne des points alignés sur une droite passant par l’origine. Pour des données discrètes, on place les points sans inventer de valeurs intermédiaires ; leur alignement avec l’origine caractérise la proportionnalité du tableau.\end{propriete}

\clearpage\rubrique{Découverte}
\begin{blocLivre}{ex-11-1}\begin{exercice}[1]{Un coefficient}
3 kg de fruits coûtent 7,50 euros, au même prix au kg. Trouver le coefficient permettant de passer des kg aux euros.
\end{exercice}\end{blocLivre}
\begin{blocLivre}{ex-11-2}\begin{exercice}[1]{Une recette doublée}
Pour quatre personnes, il faut 180 g de riz. Combien pour huit personnes ?
\end{exercice}\end{blocLivre}
\begin{blocLivre}{ex-11-3}\begin{exercice}[1]{Un pourcentage simple}
Calculer 10 \%, 50 \% et 1 \% de 240.
\end{exercice}\end{blocLivre}
\begin{blocLivre}{ex-11-4}\begin{exercice}[1]{Proportion ou effectif}
Parmi 30 élèves, 12 viennent à pied. Quel pourcentage cela représente-t-il ?
\end{exercice}\end{blocLivre}
\begin{blocLivre}{ex-11-5}\begin{exercice}[1]{Un tableau à reconnaître}
Un magasin vend des lots de stylos identiques. D’après le tableau, le prix payé est-il proportionnel au nombre de lots achetés ?\begin{center}\renewcommand{\arraystretch}{1.35}\begin{tabular}{c|ccc}\toprule
Nombre de lots & 1 & 2 & 3\\\midrule
Prix (euros) & 4 & 8 & 12\\
\bottomrule\end{tabular}\end{center}
\end{exercice}\end{blocLivre}

\rubrique{Entraînement}
\begin{blocLivre}{ex-11-6}\begin{exercice}[2]{Une recette pour six}
Une recette pour quatre demande 250 g de farine et 20 cL de lait. Donner les quantités pour six.
\end{exercice}\end{blocLivre}
\begin{blocLivre}{ex-11-7}\begin{exercice}[2]{Une remise}
Un article de 80 euros bénéficie d’une réduction de 15 \%. Calculer la réduction puis le prix payé.
\end{exercice}\end{blocLivre}
\begin{blocLivre}{ex-11-8}\begin{exercice}[2]{Une échelle de plan}
Un mur mesure 4,8 m. Quelle longueur a-t-il sur un plan à l’échelle $1:100$ ?
\label{fig-ex-11-8}\figureGeo{prop-8}
\end{exercice}\end{blocLivre}
\begin{blocLivre}{ex-11-9}\begin{exercice}[2]{Vitesse moyenne}
Un cycliste parcourt 27 km en 1 h 30. Quelle est sa vitesse moyenne en km/h ?
\end{exercice}\end{blocLivre}
\begin{blocLivre}{ex-11-10}\begin{exercice}[2]{Une situation non proportionnelle}
Une livraison coûte 4 euros fixes plus 2 euros par colis. Le prix est-il proportionnel au nombre de colis ?
\end{exercice}\end{blocLivre}

\rubrique{Problèmes}
\begin{blocLivre}{ex-11-11}\begin{exercice}[3]{Lire des points proportionnels}
Le graphique donne le prix de tickets identiques. Lire le prix de trois tickets, puis trouver le prix unitaire.
\label{fig-ex-11-11}\figureGeo{prop-11}
\end{exercice}\end{blocLivre}
\begin{blocLivre}{ex-11-12}\begin{exercice}[3]{Retrouver un effectif}
Dans un groupe, 18 personnes représentent 30 \% du total. Combien y a-t-il de personnes ?
\end{exercice}\end{blocLivre}
\begin{blocLivre}{ex-11-13}\begin{exercice}[3]{Comparer deux promotions}
Pour un article à 60 euros, A propose 20 \% de réduction ; B retire 10 euros. Quelle offre est la plus avantageuse ?
\end{exercice}\end{blocLivre}
\begin{blocLivre}{ex-11-14}\begin{exercice}[3]{Deux trajets}
On parcourt 12 km en 30 minutes puis 18 km en 1 h 30. Calculer la vitesse moyenne sur l’ensemble du trajet.
\end{exercice}\end{blocLivre}
\begin{blocLivre}{ex-11-15}\begin{exercice}[3]{Une carte et un détour}
À l’échelle $1:25000$, un itinéraire mesure 6,4 cm sur la carte. Un détour ajoute 15 \% de la distance représentée. Quelle distance réelle parcourra-t-on ?
\label{fig-ex-11-15}\figureGeo{prop-15}
\end{exercice}\end{blocLivre}

\clearpage
